% --- Archon physics formalization source begin ---
% archon:physics
% archon:covers PhyXMiniProblems/problem_phyx_mini_0509.lean
% archon:source-report reports/phyx_mini/problem_phyx_mini_0509.source.json

\chapter{Physics problem phyx\_mini\_0509}
\label{ch:PhyXMiniProblems_problem_phyx_mini_0509}

\paragraph{Problem source.}
\#\# Physical scenario

A spaceship flies past earth at a speed of \textbackslash{}( 0.990c \textbackslash{}). A crew member on board the spaceship measures its length, obtaining the value \textbackslash{}( 400 \textbackslash{}text\{ m\} \textbackslash{}) as shown in figure.

\#\# Question

What length do observers measure on earth?

\#\# Answer choices

- A: A: \textbackslash{}( 35.1 \textbackslash{}text\{ m\} \textbackslash{})

- B: B: \textbackslash{}( 27.7 \textbackslash{}text\{ m\} \textbackslash{})

- C: C: \textbackslash{}( 32.1 \textbackslash{}text\{ m\} \textbackslash{})

- D: D: \textbackslash{}( 56.4 \textbackslash{}text\{ m\} \textbackslash{})

\#\# Figure caption (auxiliary; use the image as primary evidence)

The image depicts a scene involving concepts from relativity:

1. **Rocket and Spacecraft**: On the left, there's a rocket structure with a pointed front, labeled with an initial length \textbackslash{}( l\_0 = 400 \textbackslash{}, \textbackslash{}text\{m\} \textbackslash{}) extending horizontally.

2. **Observers**: Two observers, labeled \textbackslash{}( O\_1 \textbackslash{}) and \textbackslash{}( O\_2 \textbackslash{}), sit on what appears to be Earth. Each observer holds a clock. One is drawn with distinctive spiky hair.

3. **Coordinate System**: The axes \textbackslash{}( x \textbackslash{}) and \textbackslash{}( y \textbackslash{}) are marked. The \textbackslash{}( x \textbackslash{})-axis is horizontal along the direction of the rocket, while the \textbackslash{}( y \textbackslash{})-axis is vertical.

4. **Velocity**: A green arrow pointing right indicates the rocket's velocity as \textbackslash{}( 0.990c \textbackslash{}), suggesting a speed close to the speed of light.

5. **Measurements**: Two segments are marked on the horizontal line: \textbackslash{}( x\_1 \textbackslash{}) and \textbackslash{}( x\_2 \textbackslash{}). The shorter segment near the observers is \textbackslash{}( l \textbackslash{}), presumably the contracted length due to relativistic effects.

6. **Background**: The base of the diagram includes a terrestrial surface, suggesting the observers are positioned on Earth.

The diagram illustrates the concept of length contraction in special relativity and involves observers, a moving rocket, and various notation for distances and relative velocity.

\#\# Dataset metadata

Relativity; Physical Model Grounding Reasoning, Predictive Reasoning

\paragraph{Recorded answer/context.}
D: D: \textbackslash{}( 56.4 \textbackslash{}text\{ m\} \textbackslash{})

\paragraph{Figure/image path.}
/root/proposal\_for\_physic/science-mango/phyx\_mini\_run/phyx\_data/test\_image/509.png

\paragraph{Formalization target.}
create a compiling Lean file with sorry bodies at `PhyXMiniProblems/problem\_phyx\_mini\_0509.lean`. The Lean declarations must preserve the physical quantities, dimensions or dimensional roles, figure labels, governing-law hypotheses, and final relation expressed by this problem.
Use Mathlib/Physlib names found through LeanExplore where available. If a physics API is missing, introduce faithful local abstractions rather than scalar placeholder aliases.

\begin{theorem}[Physics formalization target]
\label{thm:physics:phyx_mini_0509:target}
\lean{PhyXMiniProblems.ProblemPhyXMini0509.problem_phyx_mini_0509}
\uses{def:physics:phyx-mini-0509:phyxminiproblems-problemphyxmini0509-spaceshiplengthcontractionsetup, def:physics:phyx-mini-0509:phyxminiproblems-problemphyxmini0509-matchespassingspaceshipscenario, def:physics:phyx-mini-0509:phyxminiproblems-problemphyxmini0509-matchessuppliedspaceshipfigure, def:physics:phyx-mini-0509:phyxminiproblems-problemphyxmini0509-matchesproblemreadouts, def:physics:phyx-mini-0509:phyxminiproblems-problemphyxmini0509-hasphysicalrelativisticparameters, def:physics:phyx-mini-0509:phyxminiproblems-problemphyxmini0509-satisfiessimultaneousearthendpointmeasurement, def:physics:phyx-mini-0509:phyxminiproblems-problemphyxmini0509-satisfiesspecialrelativisticlengthcontraction, def:physics:phyx-mini-0509:phyxminiproblems-problemphyxmini0509-matchesdisplayedlengthchoice}
The assigned autoformalize agent should translate this physics problem into Lean declarations in the covered file, with theorem and lemma proof bodies written as `by sorry`.
\end{theorem}
\begin{proof}
This is an autoformalization task, not a proof task. Produce statements that can later be proved by the physics prover without weakening the physical model.
\end{proof}

% --- Archon declaration topology begin ---
\section{Declaration topology}
\begin{definition}[Length Quantity]
  \label{def:physics:phyx-mini-0509:phyxminiproblems-problemphyxmini0509-lengthquantity}
  \lean{PhyXMiniProblems.ProblemPhyXMini0509.LengthQuantity}
  A nonnegative, unit-independent physical length
\end{definition}
\begin{proof}
  This is the declared model definition of Length Quantity.
\end{proof}
\begin{definition}[length Readout]
  \label{def:physics:phyx-mini-0509:phyxminiproblems-problemphyxmini0509-lengthreadout}
  \lean{PhyXMiniProblems.ProblemPhyXMini0509.lengthReadout}
  \uses{def:physics:phyx-mini-0509:phyxminiproblems-problemphyxmini0509-lengthquantity}
  Read a physical length in a selected length unit
\end{definition}
\begin{proof}
  This is the declared model definition of length Readout.
\end{proof}
\begin{definition}[length In Meters]
  \label{def:physics:phyx-mini-0509:phyxminiproblems-problemphyxmini0509-lengthinmeters}
  \lean{PhyXMiniProblems.ProblemPhyXMini0509.lengthInMeters}
  \uses{def:physics:phyx-mini-0509:phyxminiproblems-problemphyxmini0509-lengthquantity, def:physics:phyx-mini-0509:phyxminiproblems-problemphyxmini0509-lengthreadout}
  Meter readout used by the figure and the answer choices
\end{definition}
\begin{proof}
  This is the declared model definition of length In Meters.
\end{proof}
\begin{definition}[speed In Meters Per Second]
  \label{def:physics:phyx-mini-0509:phyxminiproblems-problemphyxmini0509-speedinmeterspersecond}
  \lean{PhyXMiniProblems.ProblemPhyXMini0509.speedInMetersPerSecond}
  Read a physical speed in meters per second
\end{definition}
\begin{proof}
  This is the declared model definition of speed In Meters Per Second.
\end{proof}
\begin{definition}[vacuum Speed Of Light In Meters Per Second]
  \label{def:physics:phyx-mini-0509:phyxminiproblems-problemphyxmini0509-vacuumspeedoflightinmeterspersecond}
  \lean{PhyXMiniProblems.ProblemPhyXMini0509.vacuumSpeedOfLightInMetersPerSecond}
  Physlib's exact vacuum speed of light, read in meters per second
\end{definition}
\begin{proof}
  This is the declared model definition of vacuum Speed Of Light In Meters Per Second.
\end{proof}
\begin{definition}[Inertial Frame Label]
  \label{def:physics:phyx-mini-0509:phyxminiproblems-problemphyxmini0509-inertialframelabel}
  \lean{PhyXMiniProblems.ProblemPhyXMini0509.InertialFrameLabel}
  The two inertial-frame labels printed or implied by the figure
\end{definition}
\begin{proof}
  This is the declared model definition of Inertial Frame Label.
\end{proof}
\begin{definition}[Earth Observer]
  \label{def:physics:phyx-mini-0509:phyxminiproblems-problemphyxmini0509-earthobserver}
  \lean{PhyXMiniProblems.ProblemPhyXMini0509.EarthObserver}
  The two Earth observers shown with clocks
\end{definition}
\begin{proof}
  This is the declared model definition of Earth Observer.
\end{proof}
\begin{definition}[Spacecraft Endpoint]
  \label{def:physics:phyx-mini-0509:phyxminiproblems-problemphyxmini0509-spacecraftendpoint}
  \lean{PhyXMiniProblems.ProblemPhyXMini0509.SpacecraftEndpoint}
  The rear and front endpoints of the spaceship along its motion axis
\end{definition}
\begin{proof}
  This is the declared model definition of Spacecraft Endpoint.
\end{proof}
\begin{definition}[Position Label]
  \label{def:physics:phyx-mini-0509:phyxminiproblems-problemphyxmini0509-positionlabel}
  \lean{PhyXMiniProblems.ProblemPhyXMini0509.PositionLabel}
  Coordinate-position labels attached to the endpoint rays in the image
\end{definition}
\begin{proof}
  This is the declared model definition of Position Label.
\end{proof}
\begin{definition}[Figure Length Label]
  \label{def:physics:phyx-mini-0509:phyxminiproblems-problemphyxmini0509-figurelengthlabel}
  \lean{PhyXMiniProblems.ProblemPhyXMini0509.FigureLengthLabel}
  Length labels printed in the supplied image
\end{definition}
\begin{proof}
  This is the declared model definition of Figure Length Label.
\end{proof}
\begin{definition}[Coordinate Axis]
  \label{def:physics:phyx-mini-0509:phyxminiproblems-problemphyxmini0509-coordinateaxis}
  \lean{PhyXMiniProblems.ProblemPhyXMini0509.CoordinateAxis}
  Coordinate axes visible in the supplied schematic
\end{definition}
\begin{proof}
  This is the declared model definition of Coordinate Axis.
\end{proof}
\begin{definition}[Axis Direction]
  \label{def:physics:phyx-mini-0509:phyxminiproblems-problemphyxmini0509-axisdirection}
  \lean{PhyXMiniProblems.ProblemPhyXMini0509.AxisDirection}
  Direction of the spaceship's velocity arrow
\end{definition}
\begin{proof}
  This is the declared model definition of Axis Direction.
\end{proof}
\begin{definition}[Moving Object Kind]
  \label{def:physics:phyx-mini-0509:phyxminiproblems-problemphyxmini0509-movingobjectkind}
  \lean{PhyXMiniProblems.ProblemPhyXMini0509.MovingObjectKind}
  Physical kind of the moving object in the scenario
\end{definition}
\begin{proof}
  This is the declared model definition of Moving Object Kind.
\end{proof}
\begin{definition}[Figure Feature]
  \label{def:physics:phyx-mini-0509:phyxminiproblems-problemphyxmini0509-figurefeature}
  \lean{PhyXMiniProblems.ProblemPhyXMini0509.FigureFeature}
  Named objects and annotations directly visible in the primary image
\end{definition}
\begin{proof}
  This is the declared model definition of Figure Feature.
\end{proof}
\begin{definition}[Spaceship Length Contraction Figure]
  \label{def:physics:phyx-mini-0509:phyxminiproblems-problemphyxmini0509-spaceshiplengthcontractionfigure}
  \lean{PhyXMiniProblems.ProblemPhyXMini0509.SpaceshipLengthContractionFigure}
  \uses{def:physics:phyx-mini-0509:phyxminiproblems-problemphyxmini0509-inertialframelabel, def:physics:phyx-mini-0509:phyxminiproblems-problemphyxmini0509-earthobserver, def:physics:phyx-mini-0509:phyxminiproblems-problemphyxmini0509-spacecraftendpoint, def:physics:phyx-mini-0509:phyxminiproblems-problemphyxmini0509-positionlabel, def:physics:phyx-mini-0509:phyxminiproblems-problemphyxmini0509-figurelengthlabel, def:physics:phyx-mini-0509:phyxminiproblems-problemphyxmini0509-coordinateaxis, def:physics:phyx-mini-0509:phyxminiproblems-problemphyxmini0509-axisdirection, def:physics:phyx-mini-0509:phyxminiproblems-problemphyxmini0509-figurefeature}
  Qualitative and printed evidence from phyx\_data/test\_image/509.png. Numerical annotations are scalar readouts; the corresponding dimensionful quantities live in SpaceshipLengthContractionSetup below
\end{definition}
\begin{proof}
  This is the declared model definition of Spaceship Length Contraction Figure.
\end{proof}
\begin{definition}[Spaceship Length Contraction Setup]
  \label{def:physics:phyx-mini-0509:phyxminiproblems-problemphyxmini0509-spaceshiplengthcontractionsetup}
  \lean{PhyXMiniProblems.ProblemPhyXMini0509.SpaceshipLengthContractionSetup}
  \uses{def:physics:phyx-mini-0509:phyxminiproblems-problemphyxmini0509-lengthquantity, def:physics:phyx-mini-0509:phyxminiproblems-problemphyxmini0509-inertialframelabel, def:physics:phyx-mini-0509:phyxminiproblems-problemphyxmini0509-earthobserver, def:physics:phyx-mini-0509:phyxminiproblems-problemphyxmini0509-spacecraftendpoint, def:physics:phyx-mini-0509:phyxminiproblems-problemphyxmini0509-axisdirection, def:physics:phyx-mini-0509:phyxminiproblems-problemphyxmini0509-movingobjectkind, def:physics:phyx-mini-0509:phyxminiproblems-problemphyxmini0509-spaceshiplengthcontractionfigure}
  All independent physical quantities and coordinate readouts in the problem. crewMeasuredRestLength is measured in the spaceship rest frame. The earthMeasuredLength field is an independent Earth-frame observable; it is not defined from an answer choice or from the contraction formula
\end{definition}
\begin{proof}
  This is the declared model definition of Spaceship Length Contraction Setup.
\end{proof}
\begin{definition}[speed Fraction Of Light]
  \label{def:physics:phyx-mini-0509:phyxminiproblems-problemphyxmini0509-speedfractionoflight}
  \lean{PhyXMiniProblems.ProblemPhyXMini0509.speedFractionOfLight}
  \uses{def:physics:phyx-mini-0509:phyxminiproblems-problemphyxmini0509-speedinmeterspersecond, def:physics:phyx-mini-0509:phyxminiproblems-problemphyxmini0509-vacuumspeedoflightinmeterspersecond, def:physics:phyx-mini-0509:phyxminiproblems-problemphyxmini0509-spaceshiplengthcontractionsetup}
  The dimensionless speed parameter β = v/c
\end{definition}
\begin{proof}
  This is the declared model definition of speed Fraction Of Light.
\end{proof}
\begin{definition}[lorentz Factor]
  \label{def:physics:phyx-mini-0509:phyxminiproblems-problemphyxmini0509-lorentzfactor}
  \lean{PhyXMiniProblems.ProblemPhyXMini0509.lorentzFactor}
  \uses{def:physics:phyx-mini-0509:phyxminiproblems-problemphyxmini0509-spaceshiplengthcontractionsetup, def:physics:phyx-mini-0509:phyxminiproblems-problemphyxmini0509-speedfractionoflight}
  The Lorentz factor supplied by Physlib for the setup's v/c
\end{definition}
\begin{proof}
  This is the declared model definition of lorentz Factor.
\end{proof}
\begin{definition}[Matches Passing Spaceship Scenario]
  \label{def:physics:phyx-mini-0509:phyxminiproblems-problemphyxmini0509-matchespassingspaceshipscenario}
  \lean{PhyXMiniProblems.ProblemPhyXMini0509.MatchesPassingSpaceshipScenario}
  \uses{def:physics:phyx-mini-0509:phyxminiproblems-problemphyxmini0509-spaceshiplengthcontractionsetup}
  The crew and Earth observers use the two inertial frames in the figure
\end{definition}
\begin{proof}
  This is the declared model definition of Matches Passing Spaceship Scenario.
\end{proof}
\begin{definition}[Matches Supplied Spaceship Figure]
  \label{def:physics:phyx-mini-0509:phyxminiproblems-problemphyxmini0509-matchessuppliedspaceshipfigure}
  \lean{PhyXMiniProblems.ProblemPhyXMini0509.MatchesSuppliedSpaceshipFigure}
  \uses{def:physics:phyx-mini-0509:phyxminiproblems-problemphyxmini0509-spaceshiplengthcontractionfigure}
  Primary-image evidence: the two observers and their clocks, both coordinate axes, both frame labels, both endpoint-coordinate rays, both length labels, the rightward velocity arrow, and the 400 m and 0.990 c annotations
\end{definition}
\begin{proof}
  This is the declared model definition of Matches Supplied Spaceship Figure.
\end{proof}
\begin{definition}[Matches Problem Readouts]
  \label{def:physics:phyx-mini-0509:phyxminiproblems-problemphyxmini0509-matchesproblemreadouts}
  \lean{PhyXMiniProblems.ProblemPhyXMini0509.MatchesProblemReadouts}
  \uses{def:physics:phyx-mini-0509:phyxminiproblems-problemphyxmini0509-lengthinmeters, def:physics:phyx-mini-0509:phyxminiproblems-problemphyxmini0509-spaceshiplengthcontractionsetup, def:physics:phyx-mini-0509:phyxminiproblems-problemphyxmini0509-speedfractionoflight}
  The problem data connect the dimensionful rest length and relative speed to the two scalar annotations in the figure. No Earth-frame answer occurs here
\end{definition}
\begin{proof}
  This is the declared model definition of Matches Problem Readouts.
\end{proof}
\begin{definition}[Has Physical Relativistic Parameters]
  \label{def:physics:phyx-mini-0509:phyxminiproblems-problemphyxmini0509-hasphysicalrelativisticparameters}
  \lean{PhyXMiniProblems.ProblemPhyXMini0509.HasPhysicalRelativisticParameters}
  \uses{def:physics:phyx-mini-0509:phyxminiproblems-problemphyxmini0509-lengthinmeters, def:physics:phyx-mini-0509:phyxminiproblems-problemphyxmini0509-spaceshiplengthcontractionsetup, def:physics:phyx-mini-0509:phyxminiproblems-problemphyxmini0509-speedfractionoflight}
  Positivity, endpoint ordering, and the subluminal regime of the model
\end{definition}
\begin{proof}
  This is the declared model definition of Has Physical Relativistic Parameters.
\end{proof}
\begin{definition}[Satisfies Simultaneous Earth Endpoint Measurement]
  \label{def:physics:phyx-mini-0509:phyxminiproblems-problemphyxmini0509-satisfiessimultaneousearthendpointmeasurement}
  \lean{PhyXMiniProblems.ProblemPhyXMini0509.SatisfiesSimultaneousEarthEndpointMeasurement}
  \uses{def:physics:phyx-mini-0509:phyxminiproblems-problemphyxmini0509-lengthinmeters, def:physics:phyx-mini-0509:phyxminiproblems-problemphyxmini0509-spaceshiplengthcontractionsetup}
  Operational definition of an Earth-frame length measurement. Observers O₁ and O₂ read the rear and front positions simultaneously, and the reported length is the coordinate separation x₂ - x₁
\end{definition}
\begin{proof}
  This is the declared model definition of Satisfies Simultaneous Earth Endpoint Measurement.
\end{proof}
\begin{definition}[Satisfies Special Relativistic Length Contraction]
  \label{def:physics:phyx-mini-0509:phyxminiproblems-problemphyxmini0509-satisfiesspecialrelativisticlengthcontraction}
  \lean{PhyXMiniProblems.ProblemPhyXMini0509.SatisfiesSpecialRelativisticLengthContraction}
  \uses{def:physics:phyx-mini-0509:phyxminiproblems-problemphyxmini0509-lengthreadout, def:physics:phyx-mini-0509:phyxminiproblems-problemphyxmini0509-spaceshiplengthcontractionsetup, def:physics:phyx-mini-0509:phyxminiproblems-problemphyxmini0509-lorentzfactor}
  The generic special-relativistic length-contraction law L = L₀ / γ(β). It is stated in every common length unit and contains neither a numerical Earth-frame length nor an answer-choice label
\end{definition}
\begin{proof}
  This is the declared model definition of Satisfies Special Relativistic Length Contraction.
\end{proof}
\begin{definition}[Answer Choice]
  \label{def:physics:phyx-mini-0509:phyxminiproblems-problemphyxmini0509-answerchoice}
  \lean{PhyXMiniProblems.ProblemPhyXMini0509.AnswerChoice}
  Labels of the four meter-valued choices in the source problem
\end{definition}
\begin{proof}
  This is the declared model definition of Answer Choice.
\end{proof}
\begin{definition}[displayed Length Meters]
  \label{def:physics:phyx-mini-0509:phyxminiproblems-problemphyxmini0509-displayedlengthmeters}
  \lean{PhyXMiniProblems.ProblemPhyXMini0509.displayedLengthMeters}
  \uses{def:physics:phyx-mini-0509:phyxminiproblems-problemphyxmini0509-answerchoice}
  Displayed answer length in meters, represented exactly as a rational
\end{definition}
\begin{proof}
  This is the declared model definition of displayed Length Meters.
\end{proof}
\begin{definition}[recorded Dataset Answer]
  \label{def:physics:phyx-mini-0509:phyxminiproblems-problemphyxmini0509-recordeddatasetanswer}
  \lean{PhyXMiniProblems.ProblemPhyXMini0509.recordedDatasetAnswer}
  \uses{def:physics:phyx-mini-0509:phyxminiproblems-problemphyxmini0509-answerchoice}
  The answer label recorded by the source dataset
\end{definition}
\begin{proof}
  This is the declared model definition of recorded Dataset Answer.
\end{proof}
\begin{definition}[Matches Displayed Length Choice]
  \label{def:physics:phyx-mini-0509:phyxminiproblems-problemphyxmini0509-matchesdisplayedlengthchoice}
  \lean{PhyXMiniProblems.ProblemPhyXMini0509.MatchesDisplayedLengthChoice}
  \uses{def:physics:phyx-mini-0509:phyxminiproblems-problemphyxmini0509-lengthinmeters, def:physics:phyx-mini-0509:phyxminiproblems-problemphyxmini0509-spaceshiplengthcontractionsetup, def:physics:phyx-mini-0509:phyxminiproblems-problemphyxmini0509-answerchoice, def:physics:phyx-mini-0509:phyxminiproblems-problemphyxmini0509-displayedlengthmeters}
  Agreement with a one-decimal-place displayed answer. A tolerance of 0.05 m = 1/20 m expresses rounding to the nearest tenth of a meter
\end{definition}
\begin{proof}
  This is the declared model definition of Matches Displayed Length Choice.
\end{proof}
% --- Archon declaration topology end ---
% --- Archon physics formalization source end ---
